\documentclass[11pt,letterpaper]{article}
\usepackage[top=2.0cm, bottom=3cm, left=2.0cm, right=2.0cm]{geometry}
\usepackage[utf8]{inputenc}
\usepackage[T1]{fontenc}
\usepackage[spanish]{varioref}
\usepackage[activeacute, spanish, es-tabla]{babel}
\usepackage{fancyhdr}
\usepackage{multicol}
\usepackage{float}
\usepackage{textcomp}
\usepackage{ae,aecompl}
\usepackage{amssymb,amsmath}
\usepackage[pdftex]{graphicx}
\usepackage{askmaps}
\usepackage{multirow}
\usepackage{hyperref}
\pagestyle{fancy} 
\pagenumbering{arabic} 
\renewcommand{\headrulewidth}{0pt} 
\setlength{\headsep}{20pt} 
\setlength{\headheight}{65pt} 
\setlength{\textheight}{600pt} 
\setlength{\columnsep}{15pt}
\setlength{\parskip}{1em}
\newcommand{\universidad}{\small{Universidad Técnica Federico Santa María}}
\newcommand{\campus}{\small{Campus Santiago San Joaquín}}

% Definiciones de Título e Integrantes de Experiencia
\newcommand{\titulo}{Informe Tarea 3 \\ Arquitectura y Organización de Computadores}
\newcommand{\integrantes}{\begin{tabular}{c}
Juan Pablo Jorquera  201573533-6 \\
Cristian Navarrete 201573549-2\\
\end{tabular}}

\renewcommand{\maketitle}
{
\thispagestyle{fancy}
\begin{center}
\begin{Large}
\textbf{\titulo}\\
\end{Large}
\end{center}
\vspace{0.3cm}
}


%ENCABEZADO

\fancyhead[R]{\begin{minipage}[b]{0.405\textwidth}
\begin{center}
\universidad \\ 
\campus \\ 
\integrantes
\end{center}
\end{minipage}}
\fancyhead[L]{\vspace{15pt}\includegraphics[height=1.6cm]{Escudo.png}}
%%%%%%%%%%%%%%%%%%%%%%%%%%%%%%%%%%%%%%%%%%%%%%%%
%                                              %
% AQUI TERMINAN LAS DEFINICIONES DE ENCABEZADO %
% Y EMPIEZA EL CUERPO DEL DOCUMENTO            %
%                                              %
%%%%%%%%%%%%%%%%%%%%%%%%%%%%%%%%%%%%%%%%%%%%%%%%

\begin{document}
\setcounter{secnumdepth}{0}
\maketitle
\section{MCD}
En primer lugar se imprime para pedir los valores que se desean calcular. El problema en sí se resolvió usando el \href{https://es.wikipedia.org/wiki/Algoritmo_de_Euclides#Algoritmo_de_Euclides_tradicional}{Algoritmo de Euclides}, ya habiendo guardado los valores en $a_0$ y $a_1$ se llama a MCD como función, pasándolos a los registros $s_0$ y $s_1$, trabajando sobre ellos e intercambiándolos de ser necesario, de modo que el menor de ellos quede en $s_0$. 

Luego se creó un loop $mcd$ que va buscando el resto de la división de ambos para aplicar el Algoritmo, en caso de ser $0$ el resto significa que terminó el proceso; de no ser así, se vuelve a iterar usando como nuevos valores el menor de los antes divididos y el resto calculado.

No es necesario hacer caso base, ya que en caso de no encontrar, simplemente el resto llegará a $1$, el cual debe entregar como resto $0$.

Al ser $0$ el resto, el $beq$ dentro del loop lo dirige hacia el final, donde se imprime el valor encontrado y se almacena en la salida creada al principio. 

\section{Lucas}
Para comenzar se recibe el término enésimo como input, el cual almacenamos en $a_0$ para llamar a la función. 

Para continuar, definimos los dos primeros términos en $s_0$ y $s_1$, respectivamente y tanmbien se define $s_3$ el cual almacena la suma de los numeros vistos. A continuación se realizan saltos para entregar el resultado en caso de que los terminos $s_0$ o $s_1$ sean los términos que se busquen.

Luego se utilizó $t_0$ como un contador, el cual usaremos para verificar si se llegó al término enésimo y $t_1$ como auxiliar para ir moviendo los terminos de la suma inmediata ($L_{s-1}$ y $L_{s-2}$  )

A continuación, en el loop se guarda en variable auxiliar el termino s0 y se mueve  $s_1$ a  $s_0$ para finalmente mover el valor de  $t_1$ a  $s_1$ y realizar la suma de  $s_2$ con  $s_1$.


Al finalizar el loop se imprime el valor y se guarda en la salida definida al principio.


\section{Fibonacci}
Para comenzar se pide el input por consola para guardar el valor en $a_0$ y se llama a Fibonacci como función.

Para continuar, definimos los dos primeros términos en $s_0$ y $s_1$, respectivamente y se realizan saltos ($j$) para entregar el resultado en caso de que esos sean los términos que se busquen.

Luego se utilizó $t_0$ como un contador, el cual usaremos para verificar si se llegó al término enésimo y $t_1$ como auxiliar para ir calculando la suma de los términos actuales, para poder así actualizar $s_0$ y $s_1$ a los siguientes elementos de la serie.

Finalmente al finalizar el loop se imprime el valor y se guarda en la salida definida al principio.

Cabe destacar que como el MIPS trabaja en 32 bits y para enteros, utiliza el primer bit para el signo, el máximo término que puede calcular es cuando $n=46$, ya que es el último que cabe en $2^{31}-1$.

\section{Factorial}
Para comenzar se recibe como input el número al que se desea calcular su factorial, el cual almacenamos en $a_0$ para luego llamar a la función. 

Para continuar, definimos el primer termino  $s_0$ se realizan saltos para entregar el resultado en caso de que los terminos que se busquen sean 0 o 1.

Luego se utilizó $t_0$ como un contador, el cual usaremos para verificar si se llegó al término enésimo y ademas para realizar la multiplicacion por el siguiente termino.

Finalmente al finalizar el loop se imprime el valor y se guarda en la salida definida al principio.



\section{Extra}
Se dejo espacio en .data para ingresar la entrada y facilitar la lectura de ellos, utilizando hexadecimal para tener una mejor representación de los bits que se ingresan (comparado con decimal). 

Se dividió el trabajo en tres funciones distintas: calcular largo de bits, calcular potencia de $2$ y realizar el intercambio de bits. Desde el main se llama a la función del rotor, que utiliza las otras dos para realizar el proceso.

En primer lugar se calcula el largo de bits del primer registro, el cual almacenamos en $s_0$, el cual necesitamos para formar una constante que se usará para extraer los bits del primer registro. No es necesario hacerlo para el segundo registro, ya que los bits que se extraen en este caso se calculan desde la derecha.

Luego se busca armar una potencia de $2$ de el largo necesario para extraer los bits de cada registro: el primero se arma utilizando el largo encontrado menos la cantidad de bits que se desean intercambiar, por otro lado, la segunda potencia se arma simplemente utilizando la cantidad de bits a intercambiar, almacenándolos en $s_1$ y $s_2$ respectivamente.

Con esta potencias guardadas, se dividirán los números, para que así en $\$LO$ y en $\$HI$ (cuociente y resto) queden: la parte que se quiere incluir en el otro número y la parte que se desea guardar del número actual, dependiendo de cuál término estoy dividiendo. Estos resultados se van guardando entre $t_0$ y $t_4$.

A continuación la parte que corresponde dejar a la izquierda se amplifica por cada una de las potencias para los resultados y se les suma la parte que se desea dejar a la derecha que corresponda, quedando así los resultados en $s_3$ y $s_4$ respectivamente. Específicamente, el cuociente obtenido del primer término es lo que se va sumar al segundo término y el resto se mantiene tal cual. Por otro lado, el cuociente obtenido del segundo término se guarda, se amplifica por la potencia y se le suma lo antes calculado. Luego, el resto de éste se amplifica y luego es sumado al resto guardado del primero, para obtener así ambos términos requeridos. Dichos valores quedan en $s_3$ y $s_4$ y también se almacenan en las salidas creadas al inicio.

Es importante ver que al realizar este proceso se pueden producir problemas con los ceros a la izquierda del número, ya que son ignorados cuando se encuentran al principio del resultado, es por esto que analizaremos los casos posibles. Primero vemos que para los LSB no se generan problemas, ya que el número amplificado llena los dígitos para incluir los ceros que falten (como si se estuviera haciendo extensión de $0$ hasta tener los dígitos deseados). Por otro lado cuando haya ceros al comienzo del número en los MSB, la rutina simplemente los ignorará y guardara el resultado con menos dígitos de lo deseado. Es por esto que por ejemplo, si el resultado esperado de un número es 0b0000, en cambio se entregará 0b0, lo que corresponde hacer es obviar los $0$ al momento de leerlos.

Además, se decidió para esta ocasión no imprimir el resultado, ya que QtSpim no cuenta con impresión binaria y en su lugar lo muestra como decimal, ésto no ayuda a ver los bits que se intercambian. Entonces se recomienda verificar usando los registros $s_3$ y $s_4$ (Ya que $v_0$ se usa para salir).
\end{document}
